% Source: https://artofproblemsolving.com/community/c3h2300364_
% Stable identity: aops-topic-2300364-q01
An inequality $|x-c|<r$ says that $x$ is less than $r$ away from $c$, that is, $c-r<x<c+r$.
\begin{itemize}
\item $|x+3|<10$ means $-10<x+3<10$, i.e. $-13<x<7$ (within $10$ of $-3$).
\item $|x-5|<4$ means $-4<x-5<4$, i.e. $1<x<9$ (within $4$ of $5$).
\end{itemize}
``Or'' means $x$ may satisfy either one, so the solution set is the union of the two intervals. They overlap (on $1<x<7$), so the union is the single interval
\[\boxed{-13<x<9}.\]
On the number line, draw open circles at $-13$ and $9$ (the endpoints are not included) and shade between them:

\begin{asy}
usepackage("amsmath");
size(11cm);
defaultpen(fontsize(10pt));
draw((-16,0)--(12,0),Arrows(5));
for(int k=-15;k<=11;k+=1) draw((k,-0.15)--(k,0.15),gray);
for(int k=-15;k<=10;k+=5) label("$"+string(k)+"$",(k,-0.2),S,fontsize(7pt));
draw((-13,1.2)--(7,1.2),blue+linewidth(1.5)); draw(circle((-13,1.2),0.25),blue); draw(circle((7,1.2),0.25),blue);
label("$|x+3|<10$",(-3,1.2),N,blue+fontsize(8pt));
draw((1,2.4)--(9,2.4),red+linewidth(1.5)); draw(circle((1,2.4),0.25),red); draw(circle((9,2.4),0.25),red);
label("$|x-5|<4$",(5,2.4),N,red+fontsize(8pt));
draw((-13,0)--(9,0),linewidth(2.5)); filldraw(circle((-13,0),0.3),white,black); filldraw(circle((9,0),0.3),white,black);
draw((-13,-1.6)--(-13,-0.4),dotted); draw((9,-1.6)--(9,-0.4),dotted);
label("$-13$",(-13,-1.6),S,fontsize(8pt)); label("$9$",(9,-1.6),S,fontsize(8pt));
\end{asy}
